\section{Standard limits}

Some limits occur so often that they become standard tools. They should not be treated as isolated formulas to memorize without meaning. Each standard limit captures a basic local behavior of an important function.

\subsection*{The basic trigonometric limit}

One of the most important limits is
\[
\lim_{x\to0}\frac{\sin x}{x}=1,
\]
where \(x\) is measured in radians.

This limit says that, near zero, \(\sin x\) behaves like \(x\). The functions are not equal, but their ratio approaches \(1\).

\begin{examplebox}
Compute
\[
\lim_{x\to0}\frac{\sin 3x}{4x}.
\]
Rewrite the expression to use the standard limit:
\[
\frac{\sin 3x}{4x}=\frac34\cdot\frac{\sin 3x}{3x}.
\]
As \(x\to0\), also \(3x\to0\), so
\[
\frac{\sin 3x}{3x}\to1.
\]
Therefore
\[
\lim_{x\to0}\frac{\sin 3x}{4x}=\frac34.
\]
\end{examplebox}

The rewriting is the main step. We reshape the expression until the standard limit becomes visible.

\subsection*{Tangent and arcsine type limits}

Since
\[
\tan x=\frac{\sin x}{\cos x},
\]
and \(\cos x\to1\) as \(x\to0\), we also have
\[
\lim_{x\to0}\frac{\tan x}{x}=1.
\]
Similarly, in elementary calculus we use
\[
\lim_{x\to0}\frac{\arctan x}{x}=1.
\]
These limits express the fact that near zero, these functions are locally close to the identity function.

\begin{examplebox}
Compute
\[
\lim_{x\to0}\frac{\tan 2x}{\tan x}.
\]
Rewrite as
\[
\frac{\tan 2x}{\tan x}
=2\cdot\frac{\tan 2x}{2x}\cdot\frac{x}{\tan x}.
\]
As \(x\to0\), both fractions tend to \(1\). Therefore
\[
\lim_{x\to0}\frac{\tan 2x}{\tan x}=2.
\]
\end{examplebox}

\subsection*{The exponential standard limit}

Another standard limit is
\[
\lim_{x\to0}(1+x)^{1/x}=e.
\]
More generally, if \(u\to0\), then
\[
(1+u)^{1/u}\to e.
\]
This limit is the basis for many expressions involving variable exponents.

\begin{examplebox}
Compute
\[
\lim_{x\to0}(1-3x)^{1/x}.
\]
Let
\[
u=-3x.
\]
Then \(u\to0\) as \(x\to0\), and
\[
\frac1x=\frac{-3}{u}.
\]
Therefore
\[
(1-3x)^{1/x}=(1+u)^{-3/u}=\left((1+u)^{1/u}\right)^{-3}.
\]
Taking the limit gives
\[
e^{-3}.
\]
\end{examplebox}

This example shows why substitution is useful: it changes the expression into a recognizable standard form.

\subsection*{Standard limits as local models}

The standard limits tell us how certain functions behave near a point. Near zero,
\[
\sin x\sim x,
\qquad
\tan x\sim x,
\qquad
\arctan x\sim x.
\]
The symbol \(\sim\) means that the ratio of the two expressions tends to \(1\). This notation should be used carefully, but the idea is helpful: near zero, these functions can often be compared to simpler expressions.

\begin{summarybox}
Useful standard limits include:
\begin{itemize}
\item \(\displaystyle \lim_{x\to0}\frac{\sin x}{x}=1\),
\item \(\displaystyle \lim_{x\to0}\frac{\tan x}{x}=1\),
\item \(\displaystyle \lim_{x\to0}\frac{\arctan x}{x}=1\),
\item \(\displaystyle \lim_{x\to0}(1+x)^{1/x}=e\).
\end{itemize}
\end{summarybox}

Standard limits are not a substitute for understanding. They are reliable local models that help us simplify common limiting situations.